%%%%%%%%%%%%%%%%%%%%%%%%%%%%%%%%%%%%%%%%%
%   FEHIME CEREN AY - CURRICULUM VITAE (NORSK)
%%%%%%%%%%%%%%%%%%%%%%%%%%%%%%%%%%%%%%%%%

\documentclass[10pt,a4paper,roman]{moderncv}

\moderncvstyle{casual}
\moderncvcolor{black}

\usepackage[utf8]{inputenc}
\usepackage[norsk]{babel}
\usepackage[scale=0.75]{geometry}
\setlength{\hintscolumnwidth}{3.3cm}

%----------------------------------------------------------------------------------------
%   NAVN OG KONTAKT
%----------------------------------------------------------------------------------------

\firstname{Fehime Ceren}
\familyname{Ay}

\title{Curriculum Vitae}
\address{Bestumveien 37C, 0281}{Oslo}
\phone{+47 92 25 66 45}
\email{fcerenay@gmail.com}
\extrainfo{\href{https://cerenay.github.io}{\color{blue}cerenay.github.io}}

%----------------------------------------------------------------------------------------

\begin{document}
\makecvtitle

%----------------------------------------------------------------------------------------
%   PROFIL (KORT SAMMENDRAG)
%----------------------------------------------------------------------------------------

\section{Profil}
\cvitem{}{
Atferdsøkonom og dataanalytiker med bakgrunn fra forskning, offentlig forvaltning og næringsliv. Jeg har erfaring med å kombinere atferdsforskning, dataanalyse og eksperimentering for å skape datadrevne løsninger og bedre brukeropplevelser. Jeg trives i tverrfaglige miljøer hvor data, teknologi og innsikt møtes, og har bred erfaring med verktøy som Python, R, SQL og Databricks. Jeg har et sterkt analytisk engasjement og motiveres av å bidra til kunnskapsbaserte beslutninger.
}

%----------------------------------------------------------------------------------------
%   PERSONLIG INFORMASJON
%----------------------------------------------------------------------------------------

\section{Personlig informasjon}
\cvitem{Statsborgerskap}{Norsk, Tyrkisk}
\cvitem{Fødselsdato}{11. april 1991}
\cvitem{Nettside}{\href{https://cerenay.github.io}{cerenay.github.io}}

%----------------------------------------------------------------------------------------
%   FORSKNINGSINTERESSER
%----------------------------------------------------------------------------------------

\section{Forskningsinteresser}
\cvitem{}{Atferdsøkonomi, eksperimentell økonomi, sosial- og selvbildeeffekter, informasjon, indre og ytre motivasjon}

%----------------------------------------------------------------------------------------
%   ARBEIDSERFARING
%----------------------------------------------------------------------------------------

\section{Arbeidserfaring}

\cvitem{09/2023 -- Nåværende}{\textbf{Seniorrådgiver / Atferdsøkonom} \newline
\textit{Skatteetaten – Skattedirektoratet -  Data,Analyse og Kunnskap - Atferd og Effekt}}
\cvitem{}{
Jeg arbeider som seniorrådgiver og analytiker med prosjekter som kombinerer atferdsforskning, dataanalyse og effektmåling. Arbeidet mitt innebærer å designe og analysere eksperimenter, spørreundersøkelser og A/B-tester for å forstå og påvirke atferd i ulike deler av Skatteetaten. Jeg bruker hovedsakelig R og Python til analyse på Databricks-plattformen, og arbeider med både registerdata og surveydata.  Jeg samarbeider tett med kollegaer fra ulike divisjoner som IT, analyse og AI for å utvikle datadrevne løsninger og forbedre beslutningsgrunnlaget i etaten. Jeg deltar også i internasjonale fagmiljøer, blant annet OECDs fokusgrupper og IBPPA (International Behavioral Public Policy Association).

}

\cvitem{01/2021 -- 09/2023}{\textbf{Research Scientist / Dataanalyst – Advanced Analytics \& AI} \newline
\textit{Telenor ASA – Research \& Innovation}}
\cvitem{}{
I Advanced Analytics- og AI-teamet bidro jeg med atferdsforskning og avansert dataanalyse for Telenor ASA. Jeg jobbet med både interne og eksterne datakilder for innsamling, strukturering og analyse av data. Jeg deltok i utforming av eksperimenter og spørreundersøkelser, og brukte avanserte metoder for å analysere resultatene.

Innenfor Telenor ASA sitt virkeområde bidro jeg på tvers av forretningsenheter i Norden og Asia. Deltakelse i eksterne samarbeid med akademiske institusjoner, forskningsmiljøer og profesjonelle aktører var ogå en del av mine arbeidsoppgaver. Jeg var i tillegg tilknyttet Norges Handelshøyskole (NHH), FAIR Research Centre, som tilknyttet forsker.
}

\cvitem{08/2016 -- 12/2020}{\textbf{PhD-stipendiat} \newline
\textit{Norges Handelshøyskole (NHH) – Institutt for økonomi, FAIR}}
\cvitem{}{
Som PhD-stipendiat ved FAIR Research Centre forsket jeg på atferdsøkonomi med fokus på eksperimentelle metoder og økonomisk beslutningstaking. Arbeidet inkluderte kvantitativ analyse, eksperimentdesign og samarbeid med internasjonale forskere.
}

\cvitem{10/2014 -- 08/2016}{\textbf{Forskningsassistent} \newline
\textit{Istanbul Universitet – Fakultet for økonomi, Institutt for offentlig finans}}
\cvitem{}{
Deltok i forskningsprosjekter innen offentlig økonomi, med fokus på eksperimentell metode og økonomisk atferd.
}

%----------------------------------------------------------------------------------------
%   UTDANNING
%----------------------------------------------------------------------------------------

\section{Utdanning}

\cvitem{08/2016 -- 04/2021}{\textbf{PhD i økonomi} \newline
\textit{Norges Handelshøyskole (NHH) – FAIR, The Choice Lab} \newline
Veileder: Erik Ø. Sørensen \newline
\textit{Avhandling: "Essays on Information Preferences and Morality"} \newline
Forskningsopphold ved Brown University (jan. 2019 – juni 2019)
}

\cvitem{09/2014 -- 07/2016}{\textbf{Mastergrad i økonomi} \newline
\textit{Galatasaray Universitet, Tyrkia} \newline
\textit{Masteroppgave: "A Voting Experiment on Promises and Fairness Perceptions"}}

\cvitem{09/2009 -- 01/2014}{\textbf{Bachelorgrad i økonomi og offentlig finans} \newline
\textit{Hacettepe Universitet, Tyrkia} \newline
\textit{Bacheloroppgave: "Game Theoretical Approach to Tax Evasion and Vickrey-Clark-Groves Mechanism"}}

%----------------------------------------------------------------------------------------
%   PUBLIKASJONER
%----------------------------------------------------------------------------------------

\section{Publikasjoner (utvalg)}

\cvitem{2025}{Ay, F. C., Freddi, E., Følstad, A., Hodnebrog, S., Kvale, K., Sell, O. A., \& Ulsaker, S. (2025). \textit{Conversation logs as a source of insight: Predicting user satisfaction for customer service chatbots.} Quality and User Experience, 10(1), 1.}

\cvitem{2024}{Ay, F. C., Berge, J. W., \& Nødtvedt, K. B. (2024). \textit{Strategic curiosity: An experimental study of curiosity and dishonesty.} Journal of Economic Behavior \& Organization, 217, 287–297.}

\cvitem{2023}{Huber, C., Dreber, A., Huber, J., Johannesson, M., Kirchler, M., Weitzel, U., ... \& Holzmeister, F. (2023). \textit{Competition and moral behavior: A meta-analysis of forty-five crowd-sourced experimental designs.} PNAS, 120(23), e2215572120.}

\cvitem{2023}{Azevedo, F., Pavlovic, T., Rigo, G. G., Ay, F. C., Gjoneska, B., Etienne, T. W., ... \& Huang, G. (2023). \textit{Social and moral psychology of COVID-19 across 69 countries.} Nature Scientific Data, 10(1), 272.}

\cvitem{2023}{Brede Moe, N., Ulsaker, S., Smite, D., Moss Hildrum, J., \& Ay, F. C. (2023). \textit{Understanding the Difference between Office Presence and Co-presence in Team Member Interactions.} arXiv e-prints, arXiv-2311.}

\cvitem{2022}{Van Bavel, J. J., Cichocka, A., Capraro, V., Sjåstad, H., Nezlek, J. B., Pavlovic, T., ... \& Jørgensen, F. J. (2022). \textit{National identity predicts public health support during a global pandemic.} Nature Communications, 13(1), 517.}

%----------------------------------------------------------------------------------------
%   SPRÅK OG FERDIGHETER
%----------------------------------------------------------------------------------------

\section{Språk og ferdigheter}

\cvitem{Språk}{Norsk (flytende), Engelsk (flytende), Tyrkisk (morsmål)}
\cvitem{Programmering}{Python, R, SQL}
\cvitem{Verktøy}{Databricks, zTree, oTree, Git, Power BI}

%----------------------------------------------------------------------------------------

\end{document}